\section{Conceptos del modelo Entidad--Relación}

El modelo Entidad--Relación permite describir la estructura conceptual de un
dominio antes de decidir cómo se almacenarán físicamente sus datos. En esta
sección se distinguen sus conceptos fundamentales y las restricciones que
determinan cuándo una representación conserva correctamente la semántica del
mini--mundo \cite{elmasri2016fundamentals}.

\subsection{Tipo de relación e instancia de relación}

Un \textbf{tipo de relación} es una definición del esquema: establece qué tipos
de entidad pueden asociarse, qué papel desempeña cada participante, cuál es el
grado de la asociación y qué atributos o restricciones le corresponden. Por
ejemplo, \emph{Imparte} puede definirse entre los tipos de entidad
\emph{Profesor} y \emph{Curso}.

Una \textbf{instancia de relación}, en cambio, es una ocurrencia concreta de ese
tipo en un estado determinado de la base de datos. Si la profesora Ana Torres
imparte Bases de Datos, el par formado por esa profesora y ese curso constituye
una instancia de \emph{Imparte}. En consecuencia, el tipo pertenece al esquema
y permanece relativamente estable, mientras que sus instancias pertenecen al
estado de la base y pueden insertarse o eliminarse con el tiempo.

\begin{conceptobox}[Distinción esencial]
El tipo de relación funciona como una regla o plantilla; el conjunto de
instancias indica qué asociaciones cumplen actualmente esa regla.
\end{conceptobox}

\subsection{Migración de un atributo hacia un tipo de relación}

La migración de un atributo desde un tipo de entidad hacia una relación binaria
es adecuada cuando el valor describe realmente la \emph{asociación} entre las
entidades y no a una de ellas de manera aislada. Para que el traslado de un
atributo que sí pertenecía a una entidad sea reversible sin pérdida ni
duplicación, cada instancia de esa entidad debe participar obligatoriamente en
una sola instancia de la relación: participación total y cardinalidad máxima
igual a uno desde ese lado.

Si una entidad pudiera no participar, su valor dejaría de poder representarse
al moverlo a la relación. Si pudiera participar varias veces, el valor tendría
que repetirse en cada instancia de relación y sería necesario imponer una
restricción adicional para mantener todas las copias iguales. Cuando el valor
cambia según la contraparte, esa repetición no representa una migración
meramente estructural: revela que el atributo estaba mal ubicado y que su
dependencia semántica corresponde a la relación.

El efecto del cambio es que el atributo deja de depender exclusivamente del
identificador de la entidad y pasa a depender de la combinación de
participantes de cada instancia de relación. Por ejemplo, \emph{horasAsignadas}
no describe por sí sola a una persona ni a un proyecto, sino a la participación
de una persona en un proyecto.

\clearpage
\subsection{Tipos de relación recursiva}

Un tipo de relación es \textbf{recursivo} cuando un mismo tipo de entidad
participa más de una vez en la relación, desempeñando papeles distintos. Los
nombres de los papeles son indispensables para evitar ambigüedad
\cite{silberschatz2019database}.

Dos ejemplos son los siguientes:

\begin{itemize}
  \item En \emph{Supervisa}, el tipo de entidad \emph{Empleado} participa como
        \emph{supervisor} y como \emph{subordinado}. Un supervisor puede tener
        varios subordinados y cada subordinado puede tener, según la política
        de la organización, a lo sumo un supervisor directo.
  \item En \emph{Es prerrequisito de}, el tipo de entidad \emph{Asignatura}
        participa como \emph{prerrequisito} y como \emph{asignatura posterior}.
        Una asignatura puede requerir varias anteriores y también ser requisito
        de varias asignaturas posteriores.
\end{itemize}

\subsection{Posibilidad de emplear distintos tipos de atributos}

\subsubsection{¿Un atributo compuesto puede ser llave?}

Sí. Puede ser llave siempre que el conjunto completo de sus componentes
identifique de manera única y mínima a cada entidad. Por ejemplo, una clave
compuesta por \emph{tipoDocumento}, \emph{paísEmisor} y
\emph{númeroDocumento} puede funcionar como identificador, aunque cada
componente por separado no sea único.

\subsubsection{¿Un atributo multivaluado puede ser llave?}

No en la interpretación habitual del modelo E--R. Una llave debe proporcionar
un valor estable y unívoco por entidad, mientras que un atributo multivaluado
representa un conjunto cuyo número de elementos puede variar. Se recomienda
elegir un atributo monovaluado como identificador o transformar los valores en
un tipo de entidad relacionado que posea su propia llave.

\subsubsection{¿Un atributo derivado puede ser llave?}

No es recomendable. Aunque una expresión derivada podría ser única en un estado
particular, su valor no se almacena como hecho independiente, puede cambiar al
modificarse los datos de origen y su unicidad depende de la fórmula usada. La
llave debe formarse con atributos base, estables y controlados; el valor
derivado puede mantenerse como información calculada.

\subsubsection{¿Un atributo multivaluado puede ser compuesto?}

Sí. Cada elemento del conjunto puede tener estructura interna. Por ejemplo,
\emph{domicilios} puede ser multivaluado y cada domicilio componerse de calle,
número, colonia, código postal y municipio. La multiplicidad indica cuántos
domicilios existen; la composición describe las partes de cada uno.

\subsubsection{¿Un atributo multivaluado puede ser derivado?}

Sí, siempre que el conjunto completo pueda calcularse a partir de otros datos.
Por ejemplo, \emph{prefijosTelefónicos} puede derivarse del atributo
multivaluado \emph{teléfonos}: habrá un prefijo calculado por cada teléfono
registrado. Deben documentarse la regla de cálculo y los datos base para no
tratar el resultado derivado como una fuente independiente.

\subsubsection{¿Qué implicaría una entidad cuyos atributos fueran todos derivados?}

Implicaría que la entidad carece de estado propio almacenado. Incluso su
identificación dependería de otros hechos, por lo que no podría distinguirse ni
persistir de forma autónoma. En general, esto señala un problema de modelado:
conviene representarla como una vista o concepto derivado, como un atributo de
otra entidad, o incorporar al menos un identificador y datos base no derivados
que justifiquen su existencia independiente.

\subsection{Categorías o herencia múltiple}

Una \textbf{categoría}, también llamada tipo unión, es una subclase cuyo
conjunto de instancias procede de la unión de dos o más superclases diferentes.
Una instancia de la categoría debe pertenecer al menos a una de esas
superclases, pero no necesita pertenecer a todas. Esto la distingue de una
subclase compartida obtenida por intersección, en la que cada instancia es
simultáneamente miembro de todas las superclases involucradas
\cite{elmasri2016fundamentals}.

Ejemplos de aplicación en la vida real:

\begin{itemize}
  \item El \emph{TitularRegistral} de un inmueble puede ser una \emph{Persona}
        o una \emph{Organización}. La categoría permite relacionar ambos tipos
        con \emph{Inmueble} mediante una sola función de titularidad, conservando
        los atributos particulares de cada superclase.
  \item El \emph{BeneficiarioDePago} de una institución puede ser un
        \emph{Empleado} o un \emph{Proveedor}. La operación de pago se modela
        una sola vez sobre la categoría, mientras cada beneficiario mantiene la
        información específica del tipo al que pertenece.
\end{itemize}

La categoría evita duplicar relaciones equivalentes para cada superclase y
hace explícita la restricción de que sólo las entidades incluidas en la unión
pueden desempeñar el papel común.
